\clearpage
\section{Derivation of the Directional Derivative}
\label{app:derivations}

This appendix proves Proposition~\ref{prop:dir-deriv}.
The derivation is written for a generic softmax support.
It applies to the full vocabulary, and also to the reduced support used by the
top-$K$ construction when the tail is represented as an aggregate bucket.

\subsection{Softmax sensitivity}

We first record the first-order sensitivity of a softmax policy to a logit
perturbation, which underlies the entire derivation.

\begin{lemma}[Softmax logit sensitivity]
\label{lem:softmax-sensitivity}
Let $\tpi_i=e^{z_i}/\sum_j e^{z_j}$ be a softmax distribution with logits $z_i$, and let
$\tpi_{\eta}$ denote the distribution obtained by perturbing the logits along a fixed
direction $\bm v=(v_i)$, that is, $z_i(\eta)=z_i+\eta\,v_i$. Then
\begin{equation}
\label{eq:app-logderiv}
\frac{d}{d\eta}\log\tpi_{\eta,i}\Big|_{\eta=0}
= v_i - \sum_j \tpi_j v_j.
\end{equation}
\end{lemma}

\begin{proof}
By definition,
\begin{equation}
\log\tpi_{\eta,i}
= z_i+\eta v_i-\log\sum_j \exp(z_j+\eta v_j).
\end{equation}
Differentiating with respect to $\eta$ gives
\begin{equation}
\frac{d}{d\eta}\log\tpi_{\eta,i}
= v_i -
\frac{\sum_j v_j\exp(z_j+\eta v_j)}
{\sum_j \exp(z_j+\eta v_j)}.
\end{equation}
At $\eta=0$, the fraction is $\sum_j\tpi_j v_j$.
Thus,
\begin{equation}
\frac{d}{d\eta}\log\tpi_{\eta,i}\Big|_{\eta=0}
= v_i-\sum_j\tpi_j v_j .
\end{equation}
\end{proof}


\subsection{Directional derivative of the divergence}

\begin{proof}[Proof of Proposition~\ref{prop:dir-deriv}]
The proposition uses the unit logit direction induced by a policy-gradient step.
For the per-token surrogate $r\,\Adv$, the logit gradient is $\Adv\,r\,\nabla_z\log\tpi_k$, where
\begin{equation}
\nabla_{z_i}\log\tpi_k=\1[i=k]-\tpi_i .
\end{equation}
The scalar $r$ is positive, so the sign of the realized step is determined by
$\Adv$.
We therefore compute the derivative along the unit direction
\begin{equation}
\label{eq:app-v}
v_i = \1[i=k]-\tpi_i .
\end{equation}
Its average under the current policy is
\begin{equation}
\sum_j \tpi_j v_j
= \sum_j \tpi_j\big(\1[j=k]-\tpi_j\big)
= \tpi_k - \sum_j \tpi_j^2 .
\end{equation}

Now write the forward KL as
\begin{equation}
\KL(\bmu\|\tpi_\eta)
=
\sum_i \bmu_i\log\bmu_i
-
\sum_i \bmu_i\log\tpi_{\eta,i} .
\end{equation}
The first term is independent of $\eta$.
Applying Lemma~\ref{lem:softmax-sensitivity} to the second term gives
\begin{align}
\Taylor
&= -\sum_i \bmu_i\,
\frac{d}{d\eta}\log\tpi_{\eta,i}\Big|_{\eta=0}
\nonumber\\
&= -\sum_i \bmu_i\Big(v_i-\sum_j\tpi_j v_j\Big)
\nonumber\\
&= -\sum_i \bmu_i v_i
+ \Big(\sum_i\bmu_i\Big)\Big(\sum_j\tpi_j v_j\Big)
\nonumber\\
&= -\sum_i \bmu_i v_i + \sum_j\tpi_j v_j,
\label{eq:app-step}
\end{align}
where the last equality uses $\sum_i\bmu_i=1$.

The two terms in Eq.~(\ref{eq:app-step}) are
\begin{equation}
\sum_i \bmu_i v_i
= \sum_i \bmu_i\big(\1[i=k]-\tpi_i\big)
= \bmu_k-\sum_i\bmu_i\tpi_i,
\end{equation}
and
\begin{equation}
\sum_j\tpi_j v_j
= \tpi_k-\sum_j\tpi_j^2 .
\end{equation}
Substituting them into Eq.~(\ref{eq:app-step}) yields
\begin{align}
\Taylor
&= -\Big(\bmu_k-\sum_i\bmu_i\tpi_i\Big)
+ \Big(\tpi_k-\sum_j\tpi_j^2\Big)
\nonumber\\
&= (\tpi_k-\bmu_k)
+ \sum_i\tpi_i(\bmu_i-\tpi_i),
\label{eq:app-result}
\end{align}
which is Eq.~(\ref{eq:taylor-coeff}).
This completes the proof.
\end{proof}

\section{Equivalence under Binary-KL Approximation}
\label{app:binary-kl-equivalence}

DPPO also considers a binary KL approximation that collapses the vocabulary into
the sampled token and its complement. In this special case, the predictive
direction with the aggregated-tail estimator reduces exactly to the ratio-based
direction criterion.

Let the binary support contain the sampled token $k$ and a tail bucket. The two
policies are
\begin{equation}
\bmu^{\mathrm{bin}}=(\bmu_k, 1-\bmu_k),
\qquad
\tpi^{\mathrm{bin}}=(\tpi_k, 1-\tpi_k).
\end{equation}
The binary KL proximity criterion is
\begin{equation}
\KL\big(\bmu^{\mathrm{bin}}\|\tpi^{\mathrm{bin}}\big)
=
\bmu_k\log\frac{\bmu_k}{\tpi_k}
+
(1-\bmu_k)\log\frac{1-\bmu_k}{1-\tpi_k}.
\end{equation}
Applying Eq.~(\ref{eq:taylor-coeff}) to this two-atom support gives
\begin{align}
\Taylor^{\mathrm{bin}}
&= (\tpi_k-\bmu_k)
+ \tpi_k(\bmu_k-\tpi_k)
+ (1-\tpi_k)\big((1-\bmu_k)-(1-\tpi_k)\big)
\nonumber\\
&= (\tpi_k-\bmu_k)
- \tpi_k(\tpi_k-\bmu_k)
+ (1-\tpi_k)(\tpi_k-\bmu_k)
\nonumber\\
&= 2(1-\tpi_k)(\tpi_k-\bmu_k).
\label{eq:binary-kl-taylor}
\end{align}
For nondegenerate softmax probabilities, $0<\tpi_k<1$ and $\bmu_k>0$.
Therefore
\begin{equation}
\sign\big(\Taylor^{\mathrm{bin}}\big)
=
\sign(\tpi_k-\bmu_k)
=
\sign(r-1),
\end{equation}
where $r=\tpi_k/\bmu_k$.
It follows that
\begin{equation}
\sign\big(\Adv\,\Taylor^{\mathrm{bin}}\big)
=
\sign\big(\Adv\,(r-1)\big).
\end{equation}

Thus, if the proximity criterion is also the same binary KL, the predictive mask
with the aggregated-tail estimator is identical to DPPO-Binary-KL. This is a
degenerate case in which the binary approximation removes the distributional
degrees of freedom that the divergence-based direction criterion is designed to
use. The top-$K$ KL setting retains additional tokens, so the global correction
can differ from the sampled-token ratio direction.


\section{Predictive Divergence Masks under TV Divergence}
\label{app:topk-tv-predictive}

The paper focuses on top-$K$ forward KL because it matches the primary baseline
DPPO-TopK-KL.
DPPO can also use total variation as the proximity criterion.
This section derives the corresponding divergence-based direction criterion for TV using the same
first-order principle as in the KL case.
For a generic support, define
\begin{equation}
D_{\mathrm{TV}}(\bmu,\tpi)
=
\frac12\sum_i|\bmu_i-\tpi_i|.
\end{equation}
The divergence is not differentiable when a probability gap is exactly zero.
We use the usual subgradient convention $\sign(0)=0$.
On any fixed support, Lemma~\ref{lem:softmax-sensitivity} gives the
directional derivative of TV along the unit policy-gradient direction:
\begin{equation}
\label{eq:tv-generic-deriv}
\Taylor_{\mathrm{TV}}
=
-\frac12\sum_i s_i\,\tpi_i\Big(v_i-\sum_j\tpi_jv_j\Big),
\quad\text{where}\quad
s_i=\sign(\bmu_i-\tpi_i).
\end{equation}
The only difference from the KL case is the factor
$s_i=\sign(\bmu_i-\tpi_i)$, which comes from differentiating the absolute value
in TV.

\noindent\textbf{Aggregated-tail estimator.}
Under the top-$K$ aggregated-tail construction, the support contains the retained
tokens $\Keep$ and one tail bucket.
This treats the unseen tail as a single coordinate when computing both the TV
proximity and the softmax coupling term.
The top-$K$ TV proximity criterion is as follows:
\begin{equation}
D_{\mathrm{TV}}^{\mathrm{TopK}}(\bmu,\tpi)
=
\frac12\sum_{i\in\Keep}|\bmu_i-\tpi_i|
+
\frac12|\bmu_{\mathrm{tail}}-\tpi_{\mathrm{tail}}|.
\end{equation}
For the tail bucket, define
\begin{equation}
 s_{\mathrm{tail}}
 =\sign(\bmu_{\mathrm{tail}}-\tpi_{\mathrm{tail}}).
\end{equation}
On this reduced support,
\begin{equation}
 c_{\mathrm{agg}}
 =\sum_j\tpi_jv_j
 =\tpi_k-
 \sum_{i\in\Keep}\tpi_i^2
 -\tpi_{\mathrm{tail}}^2 .
\end{equation}
Substituting into Eq.~(\ref{eq:tv-generic-deriv}) gives
\begin{equation}
\label{eq:tv-agg-tail}
\Taylor_{\mathrm{TV}}^{\mathrm{agg}}
=
-\frac12
\left[
\sum_{i\in\Keep}
 s_i\tpi_i\big(\1[i=k]-\tpi_i-c_{\mathrm{agg}}\big)
+
 s_{\mathrm{tail}}\tpi_{\mathrm{tail}}
 \big(-\tpi_{\mathrm{tail}}-c_{\mathrm{agg}}\big)
\right].
\end{equation}

\noindent\textbf{Uniform-tail estimator.}
Under the top-$K$ uniform-tail construction, the residual mass is spread over the
$n-m$ unseen tokens, where $m=|\Keep|$ and $n$ is the vocabulary size.
This keeps the observed head terms unchanged but models the unseen tail as many
small coordinates rather than one aggregate coordinate:
\begin{equation}
\tpi_j\approx\frac{\tpi_{\mathrm{tail}}}{n-m},
\qquad
\bmu_j\approx\frac{\bmu_{\mathrm{tail}}}{n-m},
\qquad j\notin\Keep.
\end{equation}
The softmax coupling term becomes
\begin{equation}
 c_{\mathrm{uni}}
 =\tpi_k-
 \sum_{i\in\Keep}\tpi_i^2
 -\frac{\tpi_{\mathrm{tail}}^2}{n-m}.
\end{equation}
The tail contribution in Eq.~(\ref{eq:tv-generic-deriv}) then sums to
\begin{equation}
 s_{\mathrm{tail}}
 \left(
 -\frac{\tpi_{\mathrm{tail}}^2}{n-m}
 -c_{\mathrm{uni}}\tpi_{\mathrm{tail}}
 \right).
\end{equation}
Thus
\begin{equation}
\label{eq:tv-uni-tail}
\Taylor_{\mathrm{TV}}^{\mathrm{uni}}
=
-\frac12
\left[
\sum_{i\in\Keep}
 s_i\tpi_i\big(\1[i=k]-\tpi_i-c_{\mathrm{uni}}\big)
+
 s_{\mathrm{tail}}
 \left(
 -\frac{\tpi_{\mathrm{tail}}^2}{n-m}
 -c_{\mathrm{uni}}\tpi_{\mathrm{tail}}
 \right)
\right].
\end{equation}

The masking rule is unchanged.
Only the proximity measure and the divergence-based direction coefficient are replaced
by their TV counterparts.
With either estimator, the corresponding predictive TV mask is as follows:
\begin{equation}
M_t^{\mathrm{pred}} =
\begin{cases}
0, & \sign\big(\Adv_t\,\Taylor_{\mathrm{TV},t}\big)>0
\ \text{ and }\ D_{\mathrm{TV},t}^{\mathrm{TopK}}>\delta,\\[2pt]
1, & \text{otherwise.}
\end{cases}
\end{equation}
Since the paper focuses on DPPO-TopK-KL, this TV form is included as the
corresponding extension for a TV proximity criterion.
